% =============================================================================
% Circle STARK subsection for zkMetal paper
% Intended inclusion: % =============================================================================
% Circle STARK subsection for zkMetal paper
% Intended inclusion: % =============================================================================
% Circle STARK subsection for zkMetal paper
% Intended inclusion: % =============================================================================
% Circle STARK subsection for zkMetal paper
% Intended inclusion: \input{sections_circle_stark} within Section 6
% =============================================================================

\subsection{Circle STARK over Mersenne31}
\label{sec:circle-stark}

Circle STARKs~\cite{circle-stark} replace the multiplicative-subgroup FFT
domain used in classical STARKs with the \emph{circle group}
$\mathcal{C}(\mathbb{F}_p) = \{(x,y) \in \mathbb{F}_p^2 : x^2 + y^2 = 1\}$
over the Mersenne prime $p = 2^{31} - 1$.  This construction, introduced by
Hab\"ock and Levit and implemented in StarkWare's Stwo prover and
Plonky3~\cite{plonky3}, achieves two properties that are particularly
attractive for Metal GPUs: (1)~all field arithmetic fits in a single 32-bit
word, eliminating the 64-MAD Montgomery penalty of 256-bit fields; and
(2)~the circle group has order $p + 1 = 2^{31}$, providing full 31-bit
2-adicity---larger than any prime-order multiplicative subgroup at comparable
field size.

\sys implements a complete Circle STARK prover and verifier, including GPU
Circle NTT, GPU constraint evaluation, GPU Keccak-256 Merkle commitment,
and CPU FRI folding.  This subsection describes the field tower, the
proving pipeline, and the performance characteristics that emerge from
mapping this protocol to Metal's 32-bit ALU.

\subsubsection{Mersenne31 and the QM31 Extension Tower}
\label{sec:m31-tower}

The base field $\mathbb{F}_p$ with $p = 2^{31} - 1$ admits exceptionally
cheap arithmetic on 32-bit hardware.  Addition reduces to a single 32-bit
add with a conditional subtraction of $p$ (implemented branch-free via a
mask-and-shift: $(s \mathbin{\&} p) + (s \gg 31)$, where $s = a + b$).
Multiplication uses a 32$\times$32$\to$64 widening multiply followed by a
Mersenne reduction: the product $a \cdot b$ is split as
$\mathit{lo} = \mathit{prod} \mathbin{\&} p$ and
$\mathit{hi} = \mathit{prod} \gg 31$, then reduced as
$(\mathit{lo} + \mathit{hi}) \bmod p$.  This is a \emph{single MAD
instruction} on Metal's 32-bit ALU---the same cost as a BabyBear multiply
and $64\times$ cheaper than a BN254 Montgomery multiply.

Soundness of the STARK protocol requires random challenges and FRI folding
over an extension field with at least 128 bits of entropy.  Circle STARKs
use the degree-4 extension QM31, constructed as a two-step tower:
%
\begin{align}
  \text{CM31} &= \mathbb{F}_p[i] / (i^2 + 1) \label{eq:cm31} \\
  \text{QM31} &= \text{CM31}[u] / (u^2 - 2 - i) \label{eq:qm31}
\end{align}
%
CM31 elements are pairs $(a, b) \in \mathbb{F}_p^2$ representing $a + bi$,
with multiplication requiring 4 base-field multiplies (Karatsuba reduces
this to~3 in practice).  QM31 elements are pairs $(c_0, c_1) \in
\text{CM31}^2$ representing $c_0 + c_1 u$, with multiplication requiring
the identity $u^2 = 2 + i$ to reduce the cross term: $(a_0 + a_1 u)(b_0 +
b_1 u) = (a_0 b_0 + a_1 b_1 (2+i)) + (a_0 b_1 + a_1 b_0) u$.
A single QM31 multiply thus costs $3 \times 4 = 12$ base-field multiplies
(three CM31 multiplies), or equivalently 12 MAD instructions on Metal.
For comparison, a single BN254 $\mathbb{F}_r$ multiply costs $\sim$64 MADs
via CIOS---a $5.3\times$ per-operation penalty.

This cost advantage is the core thesis of small-field STARKs on Metal:
the 32-bit ALU that penalizes 256-bit fields becomes a \emph{native match}
for M31, eliminating the limb decomposition overhead entirely.

\subsubsection{Circle Group and Domain Construction}
\label{sec:circle-domain}

The circle group $\mathcal{C}(\mathbb{F}_p)$ with group operation
$(x_1, y_1) \circ (x_2, y_2) = (x_1 x_2 - y_1 y_2,\; x_1 y_2 + y_1 x_2)$
is isomorphic to the unit circle in $\mathbb{F}_p[i]$ under complex
multiplication.  The identity is $(1, 0)$ and the inverse of $(x, y)$ is
the conjugate $(x, -y)$.  The group has order $p + 1 = 2^{31}$, and \sys
uses the generator $g = (2, 1268011823)$ of order $2^{31}$.

For a trace of length $n = 2^k$, the \emph{trace domain} is the unique
order-$n$ subgroup of $\mathcal{C}(\mathbb{F}_p)$, generated by
$g^{2^{31-k}}$.  The \emph{evaluation domain} for the low-degree extension
(LDE) is a coset of a larger subgroup: with blowup factor $B$, the
evaluation domain has size $Bn$ and is obtained by multiplying the
order-$Bn$ subgroup by a fixed coset shift (the full generator $g$).  This
coset shift ensures that no evaluation point has $y = 0$, which is required
for the first FRI fold (Section~\ref{sec:circle-fri}).

\subsubsection{Proving Pipeline}
\label{sec:circle-pipeline}

The \sys Circle STARK prover follows a six-phase pipeline, with phases 1--4
executing on the GPU and phases 5--6 on the CPU:

\paragraph{Phase 1: Trace generation.}
The prover generates the execution trace as $C$ columns of $n$ M31 elements
each.  For structured AIRs (\eg Fibonacci), \sys uses a GPU witness
generation kernel that computes the trace in parallel.  For arbitrary AIRs,
the trace is generated on the CPU via the \texttt{CircleAIR} protocol and
transferred to GPU buffers via unified memory (zero-copy).

\paragraph{Phase 2: Low-degree extension via Circle NTT.}
Each trace column undergoes three steps: (a)~GPU Circle INTT on the
trace domain ($n$ points) to obtain polynomial coefficients;
(b)~zero-padding from $n$ to $Bn$ coefficients; (c)~GPU Circle NTT on
the evaluation domain ($Bn$ points) to obtain the LDE.  The Circle NTT
differs from standard radix-2 NTT in its twiddle structure: the first
butterfly layer uses $y$-coordinate twiddles (twin-coset decomposition),
while layers $1$ through $k-1$ use $x$-coordinate twiddles with the
squaring map $x \mapsto 2x^2 - 1$.  Both INTT columns and both NTT
evaluations are encoded into a single Metal command buffer to amortize
dispatch overhead.

\paragraph{Phase 3: Trace commitment.}
The LDE evaluations for each column are committed via GPU-accelerated
Keccak-256 Merkle trees.  The Merkle tree construction is fused into the
same command buffer as the NTT (Phase~2), eliminating an additional
${\sim}$0.5\,ms command buffer round-trip.  Only the 32-byte Merkle roots
are read back to the CPU for the Fiat-Shamir transcript.

\paragraph{Phase 4: Constraint evaluation and composition commitment.}
A Fiat-Shamir challenge $\alpha \in \mathbb{F}_p$ is squeezed from the
transcript.  A dedicated Metal compute kernel evaluates the AIR transition
constraints and boundary constraints at every point in the evaluation
domain, combining them into a single composition polynomial via
random linear combination with powers of $\alpha$.  The composition
polynomial is immediately committed via a Keccak Merkle tree in the
same command buffer (fused constraint evaluation + commitment).

For the Fibonacci AIR, the GPU constraint kernel evaluates two transition
constraints ($a' = b$, $b' = a + b$) and two boundary constraints
($a_0 = 1$, $b_0 = 1$) per thread.  Each thread reads two LDE values
(8~bytes), the domain $y$-coordinate (4~bytes), and writes one composition
value (4~bytes)---a memory-light, compute-light kernel that is dispatched
across the full evaluation domain.

\paragraph{Phase 5: FRI folding (CPU).}
The composition polynomial evaluations are read from the GPU buffer and
folded via the Circle FRI protocol on the CPU.  Circle FRI differs from
standard (multiplicative-group) FRI in its folding structure:

\begin{itemize}
\item \emph{First fold:} pairs points $(x, y)$ and $(x, -y)$ using the
  $y$-coordinate twiddle: $f'(x) = \tfrac{1}{2}(f(x,y) + f(x,-y)) +
  \tfrac{\alpha}{2y}(f(x,y) - f(x,-y))$.  This reduces the domain from
  the circle to its $x$-projection.
\item \emph{Subsequent folds:} pair points $x$ and $-x$ using the squaring
  map $x \mapsto 2x^2 - 1$, folding via $f'(2x^2-1) =
  \tfrac{1}{2}(f(x) + f(-x)) + \tfrac{\alpha}{2x}(f(x) - f(-x))$.
\end{itemize}

Each fold round halves the domain, commits the folded evaluations via a
Keccak Merkle tree, and squeezes a fresh challenge from the transcript.
Folding continues until the polynomial is constant.  The CPU FRI path
currently dominates proving time at larger trace sizes
(Section~\ref{sec:circle-perf}).

\paragraph{Phase 6: Query phase.}
The verifier's query indices are derived from the Fiat-Shamir transcript.
For each query index, the prover opens the trace columns and composition
polynomial at that position, providing values and Merkle authentication
paths.  Merkle paths are extracted directly from the GPU tree buffers via
pointer arithmetic on unified memory, avoiding a full tree copy to Swift
arrays.

\subsubsection{Performance Results}
\label{sec:circle-perf}

Table~\ref{tab:circle-stark} presents Circle STARK proving and verification
times across four trace sizes.  All measurements use the Fibonacci AIR
(2~columns, 2~transition constraints), a blowup factor of $4\times$
($\log_2 B = 2$), 16 FRI queries, and Keccak-256 Merkle commitments.
Timings are best-of-3 after a warmup run that primes GPU pipeline caches,
twiddle buffers, and Metal shader compilation.

\begin{table}[t]
\centering
\small
\caption{Circle STARK performance (Fibonacci AIR, $4\times$ blowup, 16 queries, M3~Pro).}
\label{tab:circle-stark}
\begin{tabular}{lrrrl}
\toprule
Trace & Prove & Verify & Proof & Notes \\
Size  & (ms)  & (ms)   & Size  & \\
\midrule
$2^{8}$  & 109  & 15  & 39\,KB & GPU dispatch overhead dominates \\
$2^{10}$ &  24  & 14  & 53\,KB & Dispatch amortized; FRI light \\
$2^{12}$ &  22  & 16  & 69\,KB & GPU phases fully utilized \\
$2^{14}$ &  21  & 28  & 87\,KB & FRI begins to dominate \\
\bottomrule
\end{tabular}
\end{table}

\paragraph{The $2^8$ anomaly.}
The most striking feature of Table~\ref{tab:circle-stark} is that $2^8$
(109\,ms) is $4.5\times$ \emph{slower} than $2^{10}$ (24\,ms) despite
processing $4\times$ less data.  This inversion is caused by Metal's
per-command-buffer dispatch overhead (${\sim}$0.5\,ms), which dominates
when the per-kernel compute is negligible.  At $2^8$, the evaluation domain
is only $4 \times 256 = 1{,}024$ elements (4\,KB), so each GPU kernel
completes in microseconds, but the prover still pays the fixed cost of
pipeline compilation, buffer allocation, and command buffer scheduling.
The warmup run amortizes shader compilation, but residual per-dispatch
overhead remains.

This establishes a \textbf{GPU dispatch floor}: below ${\sim}2^{10}$
trace size, the CPU path (\texttt{proveCPU}) would be faster, as it avoids
all Metal scheduling overhead.  \sys does not currently auto-select the CPU
path for small traces, but the crossover point is clear from the data.

\paragraph{Phase breakdown at $2^{14}$.}
Table~\ref{tab:circle-phase} profiles the proving phases at $2^{14}$ trace
size, where the evaluation domain is $2^{16} = 65{,}536$ elements.

\begin{table}[t]
\centering
\small
\caption{Circle STARK phase breakdown at $2^{14}$ trace (M3~Pro).}
\label{tab:circle-phase}
\begin{tabular}{lrr}
\toprule
Phase & Time (ms) & \% of Total \\
\midrule
Trace generation (GPU) & 0.3 & 1\% \\
NTT + trace commit (GPU) & 0.9 & 4\% \\
Fiat-Shamir & $<$0.1 & $<$1\% \\
Constraint eval + comp commit (GPU) & 2.0 & 10\% \\
FRI folding (CPU) & 12.0 & 57\% \\
Query phase & 6.0 & 28\% \\
\midrule
\textbf{Total} & \textbf{21} & 100\% \\
\bottomrule
\end{tabular}
\end{table}

FRI folding dominates at 57\% of total time, despite being executed on the
CPU.  The GPU phases (trace generation, NTT, constraint evaluation, Merkle
commitment) collectively account for only 15\% of proving time.  This
imbalance reflects two factors: (1)~the GPU phases benefit from
command buffer fusion (NTT + Merkle in one CB, constraint eval +
composition Merkle in one CB), reducing scheduling overhead to two command
buffer commits total; and (2)~CPU FRI folding involves $\log_2(Bn)$
sequential rounds of halving, each requiring a Merkle commitment and
transcript squeeze, which serializes naturally.

GPU-accelerating the FRI fold phase is a clear optimization target.  \sys
includes a \texttt{CircleFRIEngine} with GPU fold kernels (first fold via
$y$-twiddle, subsequent folds via $x$-twiddle, and a fused 2-round fold),
but the proving pipeline currently uses the CPU path for the full FRI
protocol to simplify Fiat-Shamir transcript management.  Integrating
GPU FRI folding with interleaved transcript squeezes would require
per-round GPU$\to$CPU synchronization (one command buffer commit per fold
round), which at ${\sim}$0.5\,ms per round would add ${\sim}$8\,ms of
dispatch overhead across 16~rounds---potentially negating the compute
savings at this trace size.

\paragraph{Proof sizes.}
Proof sizes range from 39\,KB ($2^8$) to 87\,KB ($2^{14}$), growing
sublinearly with trace size.  The dominant contribution is Merkle
authentication paths: each of the 16~queries opens $C = 2$ trace columns
and 1~composition polynomial, with path length $\log_2(Bn)$ hashes of
32~bytes each.  At $2^{14}$ with $4\times$ blowup, each path has 16~nodes
($2^{16}$ leaves), contributing $16 \times 3 \times 16 \times 32 =
24{,}576$~bytes per query set.  FRI round commitments and query
responses add approximately 20\,KB.

\subsubsection{Comparison with 256-bit Field Proof Systems}
\label{sec:circle-vs-plonk}

The Circle STARK's use of M31 arithmetic provides a per-operation
advantage over BN254-based systems, but STARKs inherently require more
operations than SNARKs: larger blowup factors (typically $4$--$16\times$
vs.\ none for Plonk/Groth16), more FRI rounds ($\log_2 n$ sequential
folds), and larger proofs (tens of KB vs.\ hundreds of bytes).
Table~\ref{tab:circle-vs-plonk} compares Circle STARK and Plonk at
roughly equivalent security and circuit complexity.

\begin{table}[t]
\centering
\small
\caption{Circle STARK vs.\ Plonk (BN254/KZG) on M3~Pro.}
\label{tab:circle-vs-plonk}
\begin{tabular}{lrr}
\toprule
 & Circle STARK & Plonk (KZG) \\
\midrule
Field & M31 (31-bit) & BN254 Fr (254-bit) \\
Prove (${\sim}$1K constraints) & 24\,ms & 49\,ms \\
Verify & 14\,ms & 2\,ms \\
Proof size & 53\,KB & $<$1\,KB \\
Field mul cost & 1 MAD & 64 MADs \\
FRI/PCS rounds & $\log_2 n$ & 1 (KZG) \\
Trusted setup & None & Required \\
\bottomrule
\end{tabular}
\end{table}

At roughly 1K constraints, Circle STARK proves $2\times$ faster than
Plonk---the M31 per-operation advantage ($64\times$) more than compensates
for the STARK's additional operations.  However, verification and proof
size strongly favor Plonk: KZG verification is a constant-time pairing
check (2\,ms), while STARK verification requires re-executing all query
checks and FRI verification (14\,ms).  Proof size differs by two orders of
magnitude (53\,KB vs.\ $<$1\,KB).

The crossover point depends on the application context.  For client-side
proving where proof generation latency matters most (mobile wallets,
browser provers), Circle STARK's faster proving and transparent setup are
advantages.  For on-chain verification where proof size and verification
cost dominate (ZK rollups), KZG-based systems remain preferable.

\subsubsection{Lessons for Metal GPU STARK Design}
\label{sec:circle-lessons}

The Circle STARK implementation yields several insights specific to the
Metal GPU architecture:

\paragraph{32-bit fields eliminate the ALU bottleneck.}
M31 field multiplication maps to a single 32-bit MAD instruction---Metal's
native throughput unit.  This transforms Metal's 32-bit ALU limitation
into a strength: the GPU's full datapath is utilized, unlike 256-bit
fields where $64\times$ more MAD instructions are needed per multiply.
The NTT + trace commit phase at $2^{14}$ completes in 0.9\,ms
(including two INTT, two NTT, and two Merkle trees), demonstrating that
M31 arithmetic is fast enough to make the NTT a negligible fraction of
total proving time.

\paragraph{Dispatch overhead, not compute, is the bottleneck.}
At all tested trace sizes ($2^8$ through $2^{14}$), GPU compute time for
the parallelizable phases (NTT, constraint evaluation, Merkle trees) is
under 4\,ms.  The proving time is instead dominated by sequential protocol
phases (FRI, query extraction) and Metal dispatch overhead.  This contrasts
sharply with BN254 workloads, where the 256-bit arithmetic itself dominates.
For M31 STARKs on Metal, future optimization effort should focus on
reducing the number of command buffer commits rather than accelerating
individual kernels.

\paragraph{Command buffer fusion is essential.}
\sys fuses logically distinct phases into shared command buffers: NTT +
Merkle commit in one CB, constraint evaluation + composition commit in a
second CB.  Without this fusion, the prover would require at least 6
command buffer commits (${\sim}$3\,ms of pure scheduling overhead),
increasing total time by ${\sim}$15\% at $2^{14}$.  At smaller trace
sizes, the impact is proportionally larger---explaining part of the
$2^8$ anomaly.

\paragraph{CPU FRI is a deliberate trade-off.}
Although \sys includes GPU Circle FRI fold kernels, the prover uses CPU FRI
to avoid per-round GPU$\to$CPU synchronization for Fiat-Shamir transcript
management.  At the current trace sizes ($\leq 2^{14}$), each FRI round
processes at most $2^{16}$ elements---small enough that CPU throughput
is adequate and the ${\sim}$0.5\,ms per-round dispatch overhead of a GPU
path would likely negate compute savings.  For larger traces ($2^{20}$+),
the balance shifts, and GPU FRI with batched fold rounds becomes
worthwhile.

\paragraph{The small-field thesis holds, with caveats.}
The migration from 256-bit fields to 31-bit fields delivers a clear
per-operation speedup on Metal: M31 multiply is $64\times$ cheaper than
BN254 Fr multiply, and the NTT benefits from $8\times$ better cache
utilization (4-byte vs.\ 32-byte elements).  However, the STARK protocol
structure (blowup, FRI rounds, larger proofs) partially offsets this
advantage.  The net effect at current trace sizes is a ${\sim}$2$\times$
proving speedup over Plonk at comparable constraint counts, with
significantly worse verification time and proof size.  The advantage
grows at larger trace sizes where GPU parallelism can amortize the
sequential FRI overhead.
 within Section 6
% =============================================================================

\subsection{Circle STARK over Mersenne31}
\label{sec:circle-stark}

Circle STARKs~\cite{circle-stark} replace the multiplicative-subgroup FFT
domain used in classical STARKs with the \emph{circle group}
$\mathcal{C}(\mathbb{F}_p) = \{(x,y) \in \mathbb{F}_p^2 : x^2 + y^2 = 1\}$
over the Mersenne prime $p = 2^{31} - 1$.  This construction, introduced by
Hab\"ock and Levit and implemented in StarkWare's Stwo prover and
Plonky3~\cite{plonky3}, achieves two properties that are particularly
attractive for Metal GPUs: (1)~all field arithmetic fits in a single 32-bit
word, eliminating the 64-MAD Montgomery penalty of 256-bit fields; and
(2)~the circle group has order $p + 1 = 2^{31}$, providing full 31-bit
2-adicity---larger than any prime-order multiplicative subgroup at comparable
field size.

\sys implements a complete Circle STARK prover and verifier, including GPU
Circle NTT, GPU constraint evaluation, GPU Keccak-256 Merkle commitment,
and CPU FRI folding.  This subsection describes the field tower, the
proving pipeline, and the performance characteristics that emerge from
mapping this protocol to Metal's 32-bit ALU.

\subsubsection{Mersenne31 and the QM31 Extension Tower}
\label{sec:m31-tower}

The base field $\mathbb{F}_p$ with $p = 2^{31} - 1$ admits exceptionally
cheap arithmetic on 32-bit hardware.  Addition reduces to a single 32-bit
add with a conditional subtraction of $p$ (implemented branch-free via a
mask-and-shift: $(s \mathbin{\&} p) + (s \gg 31)$, where $s = a + b$).
Multiplication uses a 32$\times$32$\to$64 widening multiply followed by a
Mersenne reduction: the product $a \cdot b$ is split as
$\mathit{lo} = \mathit{prod} \mathbin{\&} p$ and
$\mathit{hi} = \mathit{prod} \gg 31$, then reduced as
$(\mathit{lo} + \mathit{hi}) \bmod p$.  This is a \emph{single MAD
instruction} on Metal's 32-bit ALU---the same cost as a BabyBear multiply
and $64\times$ cheaper than a BN254 Montgomery multiply.

Soundness of the STARK protocol requires random challenges and FRI folding
over an extension field with at least 128 bits of entropy.  Circle STARKs
use the degree-4 extension QM31, constructed as a two-step tower:
%
\begin{align}
  \text{CM31} &= \mathbb{F}_p[i] / (i^2 + 1) \label{eq:cm31} \\
  \text{QM31} &= \text{CM31}[u] / (u^2 - 2 - i) \label{eq:qm31}
\end{align}
%
CM31 elements are pairs $(a, b) \in \mathbb{F}_p^2$ representing $a + bi$,
with multiplication requiring 4 base-field multiplies (Karatsuba reduces
this to~3 in practice).  QM31 elements are pairs $(c_0, c_1) \in
\text{CM31}^2$ representing $c_0 + c_1 u$, with multiplication requiring
the identity $u^2 = 2 + i$ to reduce the cross term: $(a_0 + a_1 u)(b_0 +
b_1 u) = (a_0 b_0 + a_1 b_1 (2+i)) + (a_0 b_1 + a_1 b_0) u$.
A single QM31 multiply thus costs $3 \times 4 = 12$ base-field multiplies
(three CM31 multiplies), or equivalently 12 MAD instructions on Metal.
For comparison, a single BN254 $\mathbb{F}_r$ multiply costs $\sim$64 MADs
via CIOS---a $5.3\times$ per-operation penalty.

This cost advantage is the core thesis of small-field STARKs on Metal:
the 32-bit ALU that penalizes 256-bit fields becomes a \emph{native match}
for M31, eliminating the limb decomposition overhead entirely.

\subsubsection{Circle Group and Domain Construction}
\label{sec:circle-domain}

The circle group $\mathcal{C}(\mathbb{F}_p)$ with group operation
$(x_1, y_1) \circ (x_2, y_2) = (x_1 x_2 - y_1 y_2,\; x_1 y_2 + y_1 x_2)$
is isomorphic to the unit circle in $\mathbb{F}_p[i]$ under complex
multiplication.  The identity is $(1, 0)$ and the inverse of $(x, y)$ is
the conjugate $(x, -y)$.  The group has order $p + 1 = 2^{31}$, and \sys
uses the generator $g = (2, 1268011823)$ of order $2^{31}$.

For a trace of length $n = 2^k$, the \emph{trace domain} is the unique
order-$n$ subgroup of $\mathcal{C}(\mathbb{F}_p)$, generated by
$g^{2^{31-k}}$.  The \emph{evaluation domain} for the low-degree extension
(LDE) is a coset of a larger subgroup: with blowup factor $B$, the
evaluation domain has size $Bn$ and is obtained by multiplying the
order-$Bn$ subgroup by a fixed coset shift (the full generator $g$).  This
coset shift ensures that no evaluation point has $y = 0$, which is required
for the first FRI fold (Section~\ref{sec:circle-fri}).

\subsubsection{Proving Pipeline}
\label{sec:circle-pipeline}

The \sys Circle STARK prover follows a six-phase pipeline, with phases 1--4
executing on the GPU and phases 5--6 on the CPU:

\paragraph{Phase 1: Trace generation.}
The prover generates the execution trace as $C$ columns of $n$ M31 elements
each.  For structured AIRs (\eg Fibonacci), \sys uses a GPU witness
generation kernel that computes the trace in parallel.  For arbitrary AIRs,
the trace is generated on the CPU via the \texttt{CircleAIR} protocol and
transferred to GPU buffers via unified memory (zero-copy).

\paragraph{Phase 2: Low-degree extension via Circle NTT.}
Each trace column undergoes three steps: (a)~GPU Circle INTT on the
trace domain ($n$ points) to obtain polynomial coefficients;
(b)~zero-padding from $n$ to $Bn$ coefficients; (c)~GPU Circle NTT on
the evaluation domain ($Bn$ points) to obtain the LDE.  The Circle NTT
differs from standard radix-2 NTT in its twiddle structure: the first
butterfly layer uses $y$-coordinate twiddles (twin-coset decomposition),
while layers $1$ through $k-1$ use $x$-coordinate twiddles with the
squaring map $x \mapsto 2x^2 - 1$.  Both INTT columns and both NTT
evaluations are encoded into a single Metal command buffer to amortize
dispatch overhead.

\paragraph{Phase 3: Trace commitment.}
The LDE evaluations for each column are committed via GPU-accelerated
Keccak-256 Merkle trees.  The Merkle tree construction is fused into the
same command buffer as the NTT (Phase~2), eliminating an additional
${\sim}$0.5\,ms command buffer round-trip.  Only the 32-byte Merkle roots
are read back to the CPU for the Fiat-Shamir transcript.

\paragraph{Phase 4: Constraint evaluation and composition commitment.}
A Fiat-Shamir challenge $\alpha \in \mathbb{F}_p$ is squeezed from the
transcript.  A dedicated Metal compute kernel evaluates the AIR transition
constraints and boundary constraints at every point in the evaluation
domain, combining them into a single composition polynomial via
random linear combination with powers of $\alpha$.  The composition
polynomial is immediately committed via a Keccak Merkle tree in the
same command buffer (fused constraint evaluation + commitment).

For the Fibonacci AIR, the GPU constraint kernel evaluates two transition
constraints ($a' = b$, $b' = a + b$) and two boundary constraints
($a_0 = 1$, $b_0 = 1$) per thread.  Each thread reads two LDE values
(8~bytes), the domain $y$-coordinate (4~bytes), and writes one composition
value (4~bytes)---a memory-light, compute-light kernel that is dispatched
across the full evaluation domain.

\paragraph{Phase 5: FRI folding (CPU).}
The composition polynomial evaluations are read from the GPU buffer and
folded via the Circle FRI protocol on the CPU.  Circle FRI differs from
standard (multiplicative-group) FRI in its folding structure:

\begin{itemize}
\item \emph{First fold:} pairs points $(x, y)$ and $(x, -y)$ using the
  $y$-coordinate twiddle: $f'(x) = \tfrac{1}{2}(f(x,y) + f(x,-y)) +
  \tfrac{\alpha}{2y}(f(x,y) - f(x,-y))$.  This reduces the domain from
  the circle to its $x$-projection.
\item \emph{Subsequent folds:} pair points $x$ and $-x$ using the squaring
  map $x \mapsto 2x^2 - 1$, folding via $f'(2x^2-1) =
  \tfrac{1}{2}(f(x) + f(-x)) + \tfrac{\alpha}{2x}(f(x) - f(-x))$.
\end{itemize}

Each fold round halves the domain, commits the folded evaluations via a
Keccak Merkle tree, and squeezes a fresh challenge from the transcript.
Folding continues until the polynomial is constant.  The CPU FRI path
currently dominates proving time at larger trace sizes
(Section~\ref{sec:circle-perf}).

\paragraph{Phase 6: Query phase.}
The verifier's query indices are derived from the Fiat-Shamir transcript.
For each query index, the prover opens the trace columns and composition
polynomial at that position, providing values and Merkle authentication
paths.  Merkle paths are extracted directly from the GPU tree buffers via
pointer arithmetic on unified memory, avoiding a full tree copy to Swift
arrays.

\subsubsection{Performance Results}
\label{sec:circle-perf}

Table~\ref{tab:circle-stark} presents Circle STARK proving and verification
times across four trace sizes.  All measurements use the Fibonacci AIR
(2~columns, 2~transition constraints), a blowup factor of $4\times$
($\log_2 B = 2$), 16 FRI queries, and Keccak-256 Merkle commitments.
Timings are best-of-3 after a warmup run that primes GPU pipeline caches,
twiddle buffers, and Metal shader compilation.

\begin{table}[t]
\centering
\small
\caption{Circle STARK performance (Fibonacci AIR, $4\times$ blowup, 16 queries, M3~Pro).}
\label{tab:circle-stark}
\begin{tabular}{lrrrl}
\toprule
Trace & Prove & Verify & Proof & Notes \\
Size  & (ms)  & (ms)   & Size  & \\
\midrule
$2^{8}$  & 109  & 15  & 39\,KB & GPU dispatch overhead dominates \\
$2^{10}$ &  24  & 14  & 53\,KB & Dispatch amortized; FRI light \\
$2^{12}$ &  22  & 16  & 69\,KB & GPU phases fully utilized \\
$2^{14}$ &  21  & 28  & 87\,KB & FRI begins to dominate \\
\bottomrule
\end{tabular}
\end{table}

\paragraph{The $2^8$ anomaly.}
The most striking feature of Table~\ref{tab:circle-stark} is that $2^8$
(109\,ms) is $4.5\times$ \emph{slower} than $2^{10}$ (24\,ms) despite
processing $4\times$ less data.  This inversion is caused by Metal's
per-command-buffer dispatch overhead (${\sim}$0.5\,ms), which dominates
when the per-kernel compute is negligible.  At $2^8$, the evaluation domain
is only $4 \times 256 = 1{,}024$ elements (4\,KB), so each GPU kernel
completes in microseconds, but the prover still pays the fixed cost of
pipeline compilation, buffer allocation, and command buffer scheduling.
The warmup run amortizes shader compilation, but residual per-dispatch
overhead remains.

This establishes a \textbf{GPU dispatch floor}: below ${\sim}2^{10}$
trace size, the CPU path (\texttt{proveCPU}) would be faster, as it avoids
all Metal scheduling overhead.  \sys does not currently auto-select the CPU
path for small traces, but the crossover point is clear from the data.

\paragraph{Phase breakdown at $2^{14}$.}
Table~\ref{tab:circle-phase} profiles the proving phases at $2^{14}$ trace
size, where the evaluation domain is $2^{16} = 65{,}536$ elements.

\begin{table}[t]
\centering
\small
\caption{Circle STARK phase breakdown at $2^{14}$ trace (M3~Pro).}
\label{tab:circle-phase}
\begin{tabular}{lrr}
\toprule
Phase & Time (ms) & \% of Total \\
\midrule
Trace generation (GPU) & 0.3 & 1\% \\
NTT + trace commit (GPU) & 0.9 & 4\% \\
Fiat-Shamir & $<$0.1 & $<$1\% \\
Constraint eval + comp commit (GPU) & 2.0 & 10\% \\
FRI folding (CPU) & 12.0 & 57\% \\
Query phase & 6.0 & 28\% \\
\midrule
\textbf{Total} & \textbf{21} & 100\% \\
\bottomrule
\end{tabular}
\end{table}

FRI folding dominates at 57\% of total time, despite being executed on the
CPU.  The GPU phases (trace generation, NTT, constraint evaluation, Merkle
commitment) collectively account for only 15\% of proving time.  This
imbalance reflects two factors: (1)~the GPU phases benefit from
command buffer fusion (NTT + Merkle in one CB, constraint eval +
composition Merkle in one CB), reducing scheduling overhead to two command
buffer commits total; and (2)~CPU FRI folding involves $\log_2(Bn)$
sequential rounds of halving, each requiring a Merkle commitment and
transcript squeeze, which serializes naturally.

GPU-accelerating the FRI fold phase is a clear optimization target.  \sys
includes a \texttt{CircleFRIEngine} with GPU fold kernels (first fold via
$y$-twiddle, subsequent folds via $x$-twiddle, and a fused 2-round fold),
but the proving pipeline currently uses the CPU path for the full FRI
protocol to simplify Fiat-Shamir transcript management.  Integrating
GPU FRI folding with interleaved transcript squeezes would require
per-round GPU$\to$CPU synchronization (one command buffer commit per fold
round), which at ${\sim}$0.5\,ms per round would add ${\sim}$8\,ms of
dispatch overhead across 16~rounds---potentially negating the compute
savings at this trace size.

\paragraph{Proof sizes.}
Proof sizes range from 39\,KB ($2^8$) to 87\,KB ($2^{14}$), growing
sublinearly with trace size.  The dominant contribution is Merkle
authentication paths: each of the 16~queries opens $C = 2$ trace columns
and 1~composition polynomial, with path length $\log_2(Bn)$ hashes of
32~bytes each.  At $2^{14}$ with $4\times$ blowup, each path has 16~nodes
($2^{16}$ leaves), contributing $16 \times 3 \times 16 \times 32 =
24{,}576$~bytes per query set.  FRI round commitments and query
responses add approximately 20\,KB.

\subsubsection{Comparison with 256-bit Field Proof Systems}
\label{sec:circle-vs-plonk}

The Circle STARK's use of M31 arithmetic provides a per-operation
advantage over BN254-based systems, but STARKs inherently require more
operations than SNARKs: larger blowup factors (typically $4$--$16\times$
vs.\ none for Plonk/Groth16), more FRI rounds ($\log_2 n$ sequential
folds), and larger proofs (tens of KB vs.\ hundreds of bytes).
Table~\ref{tab:circle-vs-plonk} compares Circle STARK and Plonk at
roughly equivalent security and circuit complexity.

\begin{table}[t]
\centering
\small
\caption{Circle STARK vs.\ Plonk (BN254/KZG) on M3~Pro.}
\label{tab:circle-vs-plonk}
\begin{tabular}{lrr}
\toprule
 & Circle STARK & Plonk (KZG) \\
\midrule
Field & M31 (31-bit) & BN254 Fr (254-bit) \\
Prove (${\sim}$1K constraints) & 24\,ms & 49\,ms \\
Verify & 14\,ms & 2\,ms \\
Proof size & 53\,KB & $<$1\,KB \\
Field mul cost & 1 MAD & 64 MADs \\
FRI/PCS rounds & $\log_2 n$ & 1 (KZG) \\
Trusted setup & None & Required \\
\bottomrule
\end{tabular}
\end{table}

At roughly 1K constraints, Circle STARK proves $2\times$ faster than
Plonk---the M31 per-operation advantage ($64\times$) more than compensates
for the STARK's additional operations.  However, verification and proof
size strongly favor Plonk: KZG verification is a constant-time pairing
check (2\,ms), while STARK verification requires re-executing all query
checks and FRI verification (14\,ms).  Proof size differs by two orders of
magnitude (53\,KB vs.\ $<$1\,KB).

The crossover point depends on the application context.  For client-side
proving where proof generation latency matters most (mobile wallets,
browser provers), Circle STARK's faster proving and transparent setup are
advantages.  For on-chain verification where proof size and verification
cost dominate (ZK rollups), KZG-based systems remain preferable.

\subsubsection{Lessons for Metal GPU STARK Design}
\label{sec:circle-lessons}

The Circle STARK implementation yields several insights specific to the
Metal GPU architecture:

\paragraph{32-bit fields eliminate the ALU bottleneck.}
M31 field multiplication maps to a single 32-bit MAD instruction---Metal's
native throughput unit.  This transforms Metal's 32-bit ALU limitation
into a strength: the GPU's full datapath is utilized, unlike 256-bit
fields where $64\times$ more MAD instructions are needed per multiply.
The NTT + trace commit phase at $2^{14}$ completes in 0.9\,ms
(including two INTT, two NTT, and two Merkle trees), demonstrating that
M31 arithmetic is fast enough to make the NTT a negligible fraction of
total proving time.

\paragraph{Dispatch overhead, not compute, is the bottleneck.}
At all tested trace sizes ($2^8$ through $2^{14}$), GPU compute time for
the parallelizable phases (NTT, constraint evaluation, Merkle trees) is
under 4\,ms.  The proving time is instead dominated by sequential protocol
phases (FRI, query extraction) and Metal dispatch overhead.  This contrasts
sharply with BN254 workloads, where the 256-bit arithmetic itself dominates.
For M31 STARKs on Metal, future optimization effort should focus on
reducing the number of command buffer commits rather than accelerating
individual kernels.

\paragraph{Command buffer fusion is essential.}
\sys fuses logically distinct phases into shared command buffers: NTT +
Merkle commit in one CB, constraint evaluation + composition commit in a
second CB.  Without this fusion, the prover would require at least 6
command buffer commits (${\sim}$3\,ms of pure scheduling overhead),
increasing total time by ${\sim}$15\% at $2^{14}$.  At smaller trace
sizes, the impact is proportionally larger---explaining part of the
$2^8$ anomaly.

\paragraph{CPU FRI is a deliberate trade-off.}
Although \sys includes GPU Circle FRI fold kernels, the prover uses CPU FRI
to avoid per-round GPU$\to$CPU synchronization for Fiat-Shamir transcript
management.  At the current trace sizes ($\leq 2^{14}$), each FRI round
processes at most $2^{16}$ elements---small enough that CPU throughput
is adequate and the ${\sim}$0.5\,ms per-round dispatch overhead of a GPU
path would likely negate compute savings.  For larger traces ($2^{20}$+),
the balance shifts, and GPU FRI with batched fold rounds becomes
worthwhile.

\paragraph{The small-field thesis holds, with caveats.}
The migration from 256-bit fields to 31-bit fields delivers a clear
per-operation speedup on Metal: M31 multiply is $64\times$ cheaper than
BN254 Fr multiply, and the NTT benefits from $8\times$ better cache
utilization (4-byte vs.\ 32-byte elements).  However, the STARK protocol
structure (blowup, FRI rounds, larger proofs) partially offsets this
advantage.  The net effect at current trace sizes is a ${\sim}$2$\times$
proving speedup over Plonk at comparable constraint counts, with
significantly worse verification time and proof size.  The advantage
grows at larger trace sizes where GPU parallelism can amortize the
sequential FRI overhead.
 within Section 6
% =============================================================================

\subsection{Circle STARK over Mersenne31}
\label{sec:circle-stark}

Circle STARKs~\cite{circle-stark} replace the multiplicative-subgroup FFT
domain used in classical STARKs with the \emph{circle group}
$\mathcal{C}(\mathbb{F}_p) = \{(x,y) \in \mathbb{F}_p^2 : x^2 + y^2 = 1\}$
over the Mersenne prime $p = 2^{31} - 1$.  This construction, introduced by
Hab\"ock and Levit and implemented in StarkWare's Stwo prover and
Plonky3~\cite{plonky3}, achieves two properties that are particularly
attractive for Metal GPUs: (1)~all field arithmetic fits in a single 32-bit
word, eliminating the 64-MAD Montgomery penalty of 256-bit fields; and
(2)~the circle group has order $p + 1 = 2^{31}$, providing full 31-bit
2-adicity---larger than any prime-order multiplicative subgroup at comparable
field size.

\sys implements a complete Circle STARK prover and verifier, including GPU
Circle NTT, GPU constraint evaluation, GPU Keccak-256 Merkle commitment,
and CPU FRI folding.  This subsection describes the field tower, the
proving pipeline, and the performance characteristics that emerge from
mapping this protocol to Metal's 32-bit ALU.

\subsubsection{Mersenne31 and the QM31 Extension Tower}
\label{sec:m31-tower}

The base field $\mathbb{F}_p$ with $p = 2^{31} - 1$ admits exceptionally
cheap arithmetic on 32-bit hardware.  Addition reduces to a single 32-bit
add with a conditional subtraction of $p$ (implemented branch-free via a
mask-and-shift: $(s \mathbin{\&} p) + (s \gg 31)$, where $s = a + b$).
Multiplication uses a 32$\times$32$\to$64 widening multiply followed by a
Mersenne reduction: the product $a \cdot b$ is split as
$\mathit{lo} = \mathit{prod} \mathbin{\&} p$ and
$\mathit{hi} = \mathit{prod} \gg 31$, then reduced as
$(\mathit{lo} + \mathit{hi}) \bmod p$.  This is a \emph{single MAD
instruction} on Metal's 32-bit ALU---the same cost as a BabyBear multiply
and $64\times$ cheaper than a BN254 Montgomery multiply.

Soundness of the STARK protocol requires random challenges and FRI folding
over an extension field with at least 128 bits of entropy.  Circle STARKs
use the degree-4 extension QM31, constructed as a two-step tower:
%
\begin{align}
  \text{CM31} &= \mathbb{F}_p[i] / (i^2 + 1) \label{eq:cm31} \\
  \text{QM31} &= \text{CM31}[u] / (u^2 - 2 - i) \label{eq:qm31}
\end{align}
%
CM31 elements are pairs $(a, b) \in \mathbb{F}_p^2$ representing $a + bi$,
with multiplication requiring 4 base-field multiplies (Karatsuba reduces
this to~3 in practice).  QM31 elements are pairs $(c_0, c_1) \in
\text{CM31}^2$ representing $c_0 + c_1 u$, with multiplication requiring
the identity $u^2 = 2 + i$ to reduce the cross term: $(a_0 + a_1 u)(b_0 +
b_1 u) = (a_0 b_0 + a_1 b_1 (2+i)) + (a_0 b_1 + a_1 b_0) u$.
A single QM31 multiply thus costs $3 \times 4 = 12$ base-field multiplies
(three CM31 multiplies), or equivalently 12 MAD instructions on Metal.
For comparison, a single BN254 $\mathbb{F}_r$ multiply costs $\sim$64 MADs
via CIOS---a $5.3\times$ per-operation penalty.

This cost advantage is the core thesis of small-field STARKs on Metal:
the 32-bit ALU that penalizes 256-bit fields becomes a \emph{native match}
for M31, eliminating the limb decomposition overhead entirely.

\subsubsection{Circle Group and Domain Construction}
\label{sec:circle-domain}

The circle group $\mathcal{C}(\mathbb{F}_p)$ with group operation
$(x_1, y_1) \circ (x_2, y_2) = (x_1 x_2 - y_1 y_2,\; x_1 y_2 + y_1 x_2)$
is isomorphic to the unit circle in $\mathbb{F}_p[i]$ under complex
multiplication.  The identity is $(1, 0)$ and the inverse of $(x, y)$ is
the conjugate $(x, -y)$.  The group has order $p + 1 = 2^{31}$, and \sys
uses the generator $g = (2, 1268011823)$ of order $2^{31}$.

For a trace of length $n = 2^k$, the \emph{trace domain} is the unique
order-$n$ subgroup of $\mathcal{C}(\mathbb{F}_p)$, generated by
$g^{2^{31-k}}$.  The \emph{evaluation domain} for the low-degree extension
(LDE) is a coset of a larger subgroup: with blowup factor $B$, the
evaluation domain has size $Bn$ and is obtained by multiplying the
order-$Bn$ subgroup by a fixed coset shift (the full generator $g$).  This
coset shift ensures that no evaluation point has $y = 0$, which is required
for the first FRI fold (Section~\ref{sec:circle-fri}).

\subsubsection{Proving Pipeline}
\label{sec:circle-pipeline}

The \sys Circle STARK prover follows a six-phase pipeline, with phases 1--4
executing on the GPU and phases 5--6 on the CPU:

\paragraph{Phase 1: Trace generation.}
The prover generates the execution trace as $C$ columns of $n$ M31 elements
each.  For structured AIRs (\eg Fibonacci), \sys uses a GPU witness
generation kernel that computes the trace in parallel.  For arbitrary AIRs,
the trace is generated on the CPU via the \texttt{CircleAIR} protocol and
transferred to GPU buffers via unified memory (zero-copy).

\paragraph{Phase 2: Low-degree extension via Circle NTT.}
Each trace column undergoes three steps: (a)~GPU Circle INTT on the
trace domain ($n$ points) to obtain polynomial coefficients;
(b)~zero-padding from $n$ to $Bn$ coefficients; (c)~GPU Circle NTT on
the evaluation domain ($Bn$ points) to obtain the LDE.  The Circle NTT
differs from standard radix-2 NTT in its twiddle structure: the first
butterfly layer uses $y$-coordinate twiddles (twin-coset decomposition),
while layers $1$ through $k-1$ use $x$-coordinate twiddles with the
squaring map $x \mapsto 2x^2 - 1$.  Both INTT columns and both NTT
evaluations are encoded into a single Metal command buffer to amortize
dispatch overhead.

\paragraph{Phase 3: Trace commitment.}
The LDE evaluations for each column are committed via GPU-accelerated
Keccak-256 Merkle trees.  The Merkle tree construction is fused into the
same command buffer as the NTT (Phase~2), eliminating an additional
${\sim}$0.5\,ms command buffer round-trip.  Only the 32-byte Merkle roots
are read back to the CPU for the Fiat-Shamir transcript.

\paragraph{Phase 4: Constraint evaluation and composition commitment.}
A Fiat-Shamir challenge $\alpha \in \mathbb{F}_p$ is squeezed from the
transcript.  A dedicated Metal compute kernel evaluates the AIR transition
constraints and boundary constraints at every point in the evaluation
domain, combining them into a single composition polynomial via
random linear combination with powers of $\alpha$.  The composition
polynomial is immediately committed via a Keccak Merkle tree in the
same command buffer (fused constraint evaluation + commitment).

For the Fibonacci AIR, the GPU constraint kernel evaluates two transition
constraints ($a' = b$, $b' = a + b$) and two boundary constraints
($a_0 = 1$, $b_0 = 1$) per thread.  Each thread reads two LDE values
(8~bytes), the domain $y$-coordinate (4~bytes), and writes one composition
value (4~bytes)---a memory-light, compute-light kernel that is dispatched
across the full evaluation domain.

\paragraph{Phase 5: FRI folding (CPU).}
The composition polynomial evaluations are read from the GPU buffer and
folded via the Circle FRI protocol on the CPU.  Circle FRI differs from
standard (multiplicative-group) FRI in its folding structure:

\begin{itemize}
\item \emph{First fold:} pairs points $(x, y)$ and $(x, -y)$ using the
  $y$-coordinate twiddle: $f'(x) = \tfrac{1}{2}(f(x,y) + f(x,-y)) +
  \tfrac{\alpha}{2y}(f(x,y) - f(x,-y))$.  This reduces the domain from
  the circle to its $x$-projection.
\item \emph{Subsequent folds:} pair points $x$ and $-x$ using the squaring
  map $x \mapsto 2x^2 - 1$, folding via $f'(2x^2-1) =
  \tfrac{1}{2}(f(x) + f(-x)) + \tfrac{\alpha}{2x}(f(x) - f(-x))$.
\end{itemize}

Each fold round halves the domain, commits the folded evaluations via a
Keccak Merkle tree, and squeezes a fresh challenge from the transcript.
Folding continues until the polynomial is constant.  The CPU FRI path
currently dominates proving time at larger trace sizes
(Section~\ref{sec:circle-perf}).

\paragraph{Phase 6: Query phase.}
The verifier's query indices are derived from the Fiat-Shamir transcript.
For each query index, the prover opens the trace columns and composition
polynomial at that position, providing values and Merkle authentication
paths.  Merkle paths are extracted directly from the GPU tree buffers via
pointer arithmetic on unified memory, avoiding a full tree copy to Swift
arrays.

\subsubsection{Performance Results}
\label{sec:circle-perf}

Table~\ref{tab:circle-stark} presents Circle STARK proving and verification
times across four trace sizes.  All measurements use the Fibonacci AIR
(2~columns, 2~transition constraints), a blowup factor of $4\times$
($\log_2 B = 2$), 16 FRI queries, and Keccak-256 Merkle commitments.
Timings are best-of-3 after a warmup run that primes GPU pipeline caches,
twiddle buffers, and Metal shader compilation.

\begin{table}[t]
\centering
\small
\caption{Circle STARK performance (Fibonacci AIR, $4\times$ blowup, 16 queries, M3~Pro).}
\label{tab:circle-stark}
\begin{tabular}{lrrrl}
\toprule
Trace & Prove & Verify & Proof & Notes \\
Size  & (ms)  & (ms)   & Size  & \\
\midrule
$2^{8}$  & 109  & 15  & 39\,KB & GPU dispatch overhead dominates \\
$2^{10}$ &  24  & 14  & 53\,KB & Dispatch amortized; FRI light \\
$2^{12}$ &  22  & 16  & 69\,KB & GPU phases fully utilized \\
$2^{14}$ &  21  & 28  & 87\,KB & FRI begins to dominate \\
\bottomrule
\end{tabular}
\end{table}

\paragraph{The $2^8$ anomaly.}
The most striking feature of Table~\ref{tab:circle-stark} is that $2^8$
(109\,ms) is $4.5\times$ \emph{slower} than $2^{10}$ (24\,ms) despite
processing $4\times$ less data.  This inversion is caused by Metal's
per-command-buffer dispatch overhead (${\sim}$0.5\,ms), which dominates
when the per-kernel compute is negligible.  At $2^8$, the evaluation domain
is only $4 \times 256 = 1{,}024$ elements (4\,KB), so each GPU kernel
completes in microseconds, but the prover still pays the fixed cost of
pipeline compilation, buffer allocation, and command buffer scheduling.
The warmup run amortizes shader compilation, but residual per-dispatch
overhead remains.

This establishes a \textbf{GPU dispatch floor}: below ${\sim}2^{10}$
trace size, the CPU path (\texttt{proveCPU}) would be faster, as it avoids
all Metal scheduling overhead.  \sys does not currently auto-select the CPU
path for small traces, but the crossover point is clear from the data.

\paragraph{Phase breakdown at $2^{14}$.}
Table~\ref{tab:circle-phase} profiles the proving phases at $2^{14}$ trace
size, where the evaluation domain is $2^{16} = 65{,}536$ elements.

\begin{table}[t]
\centering
\small
\caption{Circle STARK phase breakdown at $2^{14}$ trace (M3~Pro).}
\label{tab:circle-phase}
\begin{tabular}{lrr}
\toprule
Phase & Time (ms) & \% of Total \\
\midrule
Trace generation (GPU) & 0.3 & 1\% \\
NTT + trace commit (GPU) & 0.9 & 4\% \\
Fiat-Shamir & $<$0.1 & $<$1\% \\
Constraint eval + comp commit (GPU) & 2.0 & 10\% \\
FRI folding (CPU) & 12.0 & 57\% \\
Query phase & 6.0 & 28\% \\
\midrule
\textbf{Total} & \textbf{21} & 100\% \\
\bottomrule
\end{tabular}
\end{table}

FRI folding dominates at 57\% of total time, despite being executed on the
CPU.  The GPU phases (trace generation, NTT, constraint evaluation, Merkle
commitment) collectively account for only 15\% of proving time.  This
imbalance reflects two factors: (1)~the GPU phases benefit from
command buffer fusion (NTT + Merkle in one CB, constraint eval +
composition Merkle in one CB), reducing scheduling overhead to two command
buffer commits total; and (2)~CPU FRI folding involves $\log_2(Bn)$
sequential rounds of halving, each requiring a Merkle commitment and
transcript squeeze, which serializes naturally.

GPU-accelerating the FRI fold phase is a clear optimization target.  \sys
includes a \texttt{CircleFRIEngine} with GPU fold kernels (first fold via
$y$-twiddle, subsequent folds via $x$-twiddle, and a fused 2-round fold),
but the proving pipeline currently uses the CPU path for the full FRI
protocol to simplify Fiat-Shamir transcript management.  Integrating
GPU FRI folding with interleaved transcript squeezes would require
per-round GPU$\to$CPU synchronization (one command buffer commit per fold
round), which at ${\sim}$0.5\,ms per round would add ${\sim}$8\,ms of
dispatch overhead across 16~rounds---potentially negating the compute
savings at this trace size.

\paragraph{Proof sizes.}
Proof sizes range from 39\,KB ($2^8$) to 87\,KB ($2^{14}$), growing
sublinearly with trace size.  The dominant contribution is Merkle
authentication paths: each of the 16~queries opens $C = 2$ trace columns
and 1~composition polynomial, with path length $\log_2(Bn)$ hashes of
32~bytes each.  At $2^{14}$ with $4\times$ blowup, each path has 16~nodes
($2^{16}$ leaves), contributing $16 \times 3 \times 16 \times 32 =
24{,}576$~bytes per query set.  FRI round commitments and query
responses add approximately 20\,KB.

\subsubsection{Comparison with 256-bit Field Proof Systems}
\label{sec:circle-vs-plonk}

The Circle STARK's use of M31 arithmetic provides a per-operation
advantage over BN254-based systems, but STARKs inherently require more
operations than SNARKs: larger blowup factors (typically $4$--$16\times$
vs.\ none for Plonk/Groth16), more FRI rounds ($\log_2 n$ sequential
folds), and larger proofs (tens of KB vs.\ hundreds of bytes).
Table~\ref{tab:circle-vs-plonk} compares Circle STARK and Plonk at
roughly equivalent security and circuit complexity.

\begin{table}[t]
\centering
\small
\caption{Circle STARK vs.\ Plonk (BN254/KZG) on M3~Pro.}
\label{tab:circle-vs-plonk}
\begin{tabular}{lrr}
\toprule
 & Circle STARK & Plonk (KZG) \\
\midrule
Field & M31 (31-bit) & BN254 Fr (254-bit) \\
Prove (${\sim}$1K constraints) & 24\,ms & 49\,ms \\
Verify & 14\,ms & 2\,ms \\
Proof size & 53\,KB & $<$1\,KB \\
Field mul cost & 1 MAD & 64 MADs \\
FRI/PCS rounds & $\log_2 n$ & 1 (KZG) \\
Trusted setup & None & Required \\
\bottomrule
\end{tabular}
\end{table}

At roughly 1K constraints, Circle STARK proves $2\times$ faster than
Plonk---the M31 per-operation advantage ($64\times$) more than compensates
for the STARK's additional operations.  However, verification and proof
size strongly favor Plonk: KZG verification is a constant-time pairing
check (2\,ms), while STARK verification requires re-executing all query
checks and FRI verification (14\,ms).  Proof size differs by two orders of
magnitude (53\,KB vs.\ $<$1\,KB).

The crossover point depends on the application context.  For client-side
proving where proof generation latency matters most (mobile wallets,
browser provers), Circle STARK's faster proving and transparent setup are
advantages.  For on-chain verification where proof size and verification
cost dominate (ZK rollups), KZG-based systems remain preferable.

\subsubsection{Lessons for Metal GPU STARK Design}
\label{sec:circle-lessons}

The Circle STARK implementation yields several insights specific to the
Metal GPU architecture:

\paragraph{32-bit fields eliminate the ALU bottleneck.}
M31 field multiplication maps to a single 32-bit MAD instruction---Metal's
native throughput unit.  This transforms Metal's 32-bit ALU limitation
into a strength: the GPU's full datapath is utilized, unlike 256-bit
fields where $64\times$ more MAD instructions are needed per multiply.
The NTT + trace commit phase at $2^{14}$ completes in 0.9\,ms
(including two INTT, two NTT, and two Merkle trees), demonstrating that
M31 arithmetic is fast enough to make the NTT a negligible fraction of
total proving time.

\paragraph{Dispatch overhead, not compute, is the bottleneck.}
At all tested trace sizes ($2^8$ through $2^{14}$), GPU compute time for
the parallelizable phases (NTT, constraint evaluation, Merkle trees) is
under 4\,ms.  The proving time is instead dominated by sequential protocol
phases (FRI, query extraction) and Metal dispatch overhead.  This contrasts
sharply with BN254 workloads, where the 256-bit arithmetic itself dominates.
For M31 STARKs on Metal, future optimization effort should focus on
reducing the number of command buffer commits rather than accelerating
individual kernels.

\paragraph{Command buffer fusion is essential.}
\sys fuses logically distinct phases into shared command buffers: NTT +
Merkle commit in one CB, constraint evaluation + composition commit in a
second CB.  Without this fusion, the prover would require at least 6
command buffer commits (${\sim}$3\,ms of pure scheduling overhead),
increasing total time by ${\sim}$15\% at $2^{14}$.  At smaller trace
sizes, the impact is proportionally larger---explaining part of the
$2^8$ anomaly.

\paragraph{CPU FRI is a deliberate trade-off.}
Although \sys includes GPU Circle FRI fold kernels, the prover uses CPU FRI
to avoid per-round GPU$\to$CPU synchronization for Fiat-Shamir transcript
management.  At the current trace sizes ($\leq 2^{14}$), each FRI round
processes at most $2^{16}$ elements---small enough that CPU throughput
is adequate and the ${\sim}$0.5\,ms per-round dispatch overhead of a GPU
path would likely negate compute savings.  For larger traces ($2^{20}$+),
the balance shifts, and GPU FRI with batched fold rounds becomes
worthwhile.

\paragraph{The small-field thesis holds, with caveats.}
The migration from 256-bit fields to 31-bit fields delivers a clear
per-operation speedup on Metal: M31 multiply is $64\times$ cheaper than
BN254 Fr multiply, and the NTT benefits from $8\times$ better cache
utilization (4-byte vs.\ 32-byte elements).  However, the STARK protocol
structure (blowup, FRI rounds, larger proofs) partially offsets this
advantage.  The net effect at current trace sizes is a ${\sim}$2$\times$
proving speedup over Plonk at comparable constraint counts, with
significantly worse verification time and proof size.  The advantage
grows at larger trace sizes where GPU parallelism can amortize the
sequential FRI overhead.
 within Section 6
% =============================================================================

\subsection{Circle STARK over Mersenne31}
\label{sec:circle-stark}

Circle STARKs~\cite{circle-stark} replace the multiplicative-subgroup FFT
domain used in classical STARKs with the \emph{circle group}
$\mathcal{C}(\mathbb{F}_p) = \{(x,y) \in \mathbb{F}_p^2 : x^2 + y^2 = 1\}$
over the Mersenne prime $p = 2^{31} - 1$.  This construction, introduced by
Hab\"ock and Levit and implemented in StarkWare's Stwo prover and
Plonky3~\cite{plonky3}, achieves two properties that are particularly
attractive for Metal GPUs: (1)~all field arithmetic fits in a single 32-bit
word, eliminating the 64-MAD Montgomery penalty of 256-bit fields; and
(2)~the circle group has order $p + 1 = 2^{31}$, providing full 31-bit
2-adicity---larger than any prime-order multiplicative subgroup at comparable
field size.

\sys implements a complete Circle STARK prover and verifier, including GPU
Circle NTT, GPU constraint evaluation, GPU Keccak-256 Merkle commitment,
and CPU FRI folding.  This subsection describes the field tower, the
proving pipeline, and the performance characteristics that emerge from
mapping this protocol to Metal's 32-bit ALU.

\subsubsection{Mersenne31 and the QM31 Extension Tower}
\label{sec:m31-tower}

The base field $\mathbb{F}_p$ with $p = 2^{31} - 1$ admits exceptionally
cheap arithmetic on 32-bit hardware.  Addition reduces to a single 32-bit
add with a conditional subtraction of $p$ (implemented branch-free via a
mask-and-shift: $(s \mathbin{\&} p) + (s \gg 31)$, where $s = a + b$).
Multiplication uses a 32$\times$32$\to$64 widening multiply followed by a
Mersenne reduction: the product $a \cdot b$ is split as
$\mathit{lo} = \mathit{prod} \mathbin{\&} p$ and
$\mathit{hi} = \mathit{prod} \gg 31$, then reduced as
$(\mathit{lo} + \mathit{hi}) \bmod p$.  This is a \emph{single MAD
instruction} on Metal's 32-bit ALU---the same cost as a BabyBear multiply
and $64\times$ cheaper than a BN254 Montgomery multiply.

Soundness of the STARK protocol requires random challenges and FRI folding
over an extension field with at least 128 bits of entropy.  Circle STARKs
use the degree-4 extension QM31, constructed as a two-step tower:
%
\begin{align}
  \text{CM31} &= \mathbb{F}_p[i] / (i^2 + 1) \label{eq:cm31} \\
  \text{QM31} &= \text{CM31}[u] / (u^2 - 2 - i) \label{eq:qm31}
\end{align}
%
CM31 elements are pairs $(a, b) \in \mathbb{F}_p^2$ representing $a + bi$,
with multiplication requiring 4 base-field multiplies (Karatsuba reduces
this to~3 in practice).  QM31 elements are pairs $(c_0, c_1) \in
\text{CM31}^2$ representing $c_0 + c_1 u$, with multiplication requiring
the identity $u^2 = 2 + i$ to reduce the cross term: $(a_0 + a_1 u)(b_0 +
b_1 u) = (a_0 b_0 + a_1 b_1 (2+i)) + (a_0 b_1 + a_1 b_0) u$.
A single QM31 multiply thus costs $3 \times 4 = 12$ base-field multiplies
(three CM31 multiplies), or equivalently 12 MAD instructions on Metal.
For comparison, a single BN254 $\mathbb{F}_r$ multiply costs $\sim$64 MADs
via CIOS---a $5.3\times$ per-operation penalty.

This cost advantage is the core thesis of small-field STARKs on Metal:
the 32-bit ALU that penalizes 256-bit fields becomes a \emph{native match}
for M31, eliminating the limb decomposition overhead entirely.

\subsubsection{Circle Group and Domain Construction}
\label{sec:circle-domain}

The circle group $\mathcal{C}(\mathbb{F}_p)$ with group operation
$(x_1, y_1) \circ (x_2, y_2) = (x_1 x_2 - y_1 y_2,\; x_1 y_2 + y_1 x_2)$
is isomorphic to the unit circle in $\mathbb{F}_p[i]$ under complex
multiplication.  The identity is $(1, 0)$ and the inverse of $(x, y)$ is
the conjugate $(x, -y)$.  The group has order $p + 1 = 2^{31}$, and \sys
uses the generator $g = (2, 1268011823)$ of order $2^{31}$.

For a trace of length $n = 2^k$, the \emph{trace domain} is the unique
order-$n$ subgroup of $\mathcal{C}(\mathbb{F}_p)$, generated by
$g^{2^{31-k}}$.  The \emph{evaluation domain} for the low-degree extension
(LDE) is a coset of a larger subgroup: with blowup factor $B$, the
evaluation domain has size $Bn$ and is obtained by multiplying the
order-$Bn$ subgroup by a fixed coset shift (the full generator $g$).  This
coset shift ensures that no evaluation point has $y = 0$, which is required
for the first FRI fold (Section~\ref{sec:circle-fri}).

\subsubsection{Proving Pipeline}
\label{sec:circle-pipeline}

The \sys Circle STARK prover follows a six-phase pipeline, with phases 1--4
executing on the GPU and phases 5--6 on the CPU:

\paragraph{Phase 1: Trace generation.}
The prover generates the execution trace as $C$ columns of $n$ M31 elements
each.  For structured AIRs (\eg Fibonacci), \sys uses a GPU witness
generation kernel that computes the trace in parallel.  For arbitrary AIRs,
the trace is generated on the CPU via the \texttt{CircleAIR} protocol and
transferred to GPU buffers via unified memory (zero-copy).

\paragraph{Phase 2: Low-degree extension via Circle NTT.}
Each trace column undergoes three steps: (a)~GPU Circle INTT on the
trace domain ($n$ points) to obtain polynomial coefficients;
(b)~zero-padding from $n$ to $Bn$ coefficients; (c)~GPU Circle NTT on
the evaluation domain ($Bn$ points) to obtain the LDE.  The Circle NTT
differs from standard radix-2 NTT in its twiddle structure: the first
butterfly layer uses $y$-coordinate twiddles (twin-coset decomposition),
while layers $1$ through $k-1$ use $x$-coordinate twiddles with the
squaring map $x \mapsto 2x^2 - 1$.  Both INTT columns and both NTT
evaluations are encoded into a single Metal command buffer to amortize
dispatch overhead.

\paragraph{Phase 3: Trace commitment.}
The LDE evaluations for each column are committed via GPU-accelerated
Keccak-256 Merkle trees.  The Merkle tree construction is fused into the
same command buffer as the NTT (Phase~2), eliminating an additional
${\sim}$0.5\,ms command buffer round-trip.  Only the 32-byte Merkle roots
are read back to the CPU for the Fiat-Shamir transcript.

\paragraph{Phase 4: Constraint evaluation and composition commitment.}
A Fiat-Shamir challenge $\alpha \in \mathbb{F}_p$ is squeezed from the
transcript.  A dedicated Metal compute kernel evaluates the AIR transition
constraints and boundary constraints at every point in the evaluation
domain, combining them into a single composition polynomial via
random linear combination with powers of $\alpha$.  The composition
polynomial is immediately committed via a Keccak Merkle tree in the
same command buffer (fused constraint evaluation + commitment).

For the Fibonacci AIR, the GPU constraint kernel evaluates two transition
constraints ($a' = b$, $b' = a + b$) and two boundary constraints
($a_0 = 1$, $b_0 = 1$) per thread.  Each thread reads two LDE values
(8~bytes), the domain $y$-coordinate (4~bytes), and writes one composition
value (4~bytes)---a memory-light, compute-light kernel that is dispatched
across the full evaluation domain.

\paragraph{Phase 5: FRI folding (CPU).}
The composition polynomial evaluations are read from the GPU buffer and
folded via the Circle FRI protocol on the CPU.  Circle FRI differs from
standard (multiplicative-group) FRI in its folding structure:

\begin{itemize}
\item \emph{First fold:} pairs points $(x, y)$ and $(x, -y)$ using the
  $y$-coordinate twiddle: $f'(x) = \tfrac{1}{2}(f(x,y) + f(x,-y)) +
  \tfrac{\alpha}{2y}(f(x,y) - f(x,-y))$.  This reduces the domain from
  the circle to its $x$-projection.
\item \emph{Subsequent folds:} pair points $x$ and $-x$ using the squaring
  map $x \mapsto 2x^2 - 1$, folding via $f'(2x^2-1) =
  \tfrac{1}{2}(f(x) + f(-x)) + \tfrac{\alpha}{2x}(f(x) - f(-x))$.
\end{itemize}

Each fold round halves the domain, commits the folded evaluations via a
Keccak Merkle tree, and squeezes a fresh challenge from the transcript.
Folding continues until the polynomial is constant.  The CPU FRI path
currently dominates proving time at larger trace sizes
(Section~\ref{sec:circle-perf}).

\paragraph{Phase 6: Query phase.}
The verifier's query indices are derived from the Fiat-Shamir transcript.
For each query index, the prover opens the trace columns and composition
polynomial at that position, providing values and Merkle authentication
paths.  Merkle paths are extracted directly from the GPU tree buffers via
pointer arithmetic on unified memory, avoiding a full tree copy to Swift
arrays.

\subsubsection{Performance Results}
\label{sec:circle-perf}

Table~\ref{tab:circle-stark} presents Circle STARK proving and verification
times across four trace sizes.  All measurements use the Fibonacci AIR
(2~columns, 2~transition constraints), a blowup factor of $4\times$
($\log_2 B = 2$), 16 FRI queries, and Keccak-256 Merkle commitments.
Timings are best-of-3 after a warmup run that primes GPU pipeline caches,
twiddle buffers, and Metal shader compilation.

\begin{table}[t]
\centering
\small
\caption{Circle STARK performance (Fibonacci AIR, $4\times$ blowup, 16 queries, M3~Pro).}
\label{tab:circle-stark}
\begin{tabular}{lrrrl}
\toprule
Trace & Prove & Verify & Proof & Notes \\
Size  & (ms)  & (ms)   & Size  & \\
\midrule
$2^{8}$  & 109  & 15  & 39\,KB & GPU dispatch overhead dominates \\
$2^{10}$ &  24  & 14  & 53\,KB & Dispatch amortized; FRI light \\
$2^{12}$ &  22  & 16  & 69\,KB & GPU phases fully utilized \\
$2^{14}$ &  21  & 28  & 87\,KB & FRI begins to dominate \\
\bottomrule
\end{tabular}
\end{table}

\paragraph{The $2^8$ anomaly.}
The most striking feature of Table~\ref{tab:circle-stark} is that $2^8$
(109\,ms) is $4.5\times$ \emph{slower} than $2^{10}$ (24\,ms) despite
processing $4\times$ less data.  This inversion is caused by Metal's
per-command-buffer dispatch overhead (${\sim}$0.5\,ms), which dominates
when the per-kernel compute is negligible.  At $2^8$, the evaluation domain
is only $4 \times 256 = 1{,}024$ elements (4\,KB), so each GPU kernel
completes in microseconds, but the prover still pays the fixed cost of
pipeline compilation, buffer allocation, and command buffer scheduling.
The warmup run amortizes shader compilation, but residual per-dispatch
overhead remains.

This establishes a \textbf{GPU dispatch floor}: below ${\sim}2^{10}$
trace size, the CPU path (\texttt{proveCPU}) would be faster, as it avoids
all Metal scheduling overhead.  \sys does not currently auto-select the CPU
path for small traces, but the crossover point is clear from the data.

\paragraph{Phase breakdown at $2^{14}$.}
Table~\ref{tab:circle-phase} profiles the proving phases at $2^{14}$ trace
size, where the evaluation domain is $2^{16} = 65{,}536$ elements.

\begin{table}[t]
\centering
\small
\caption{Circle STARK phase breakdown at $2^{14}$ trace (M3~Pro).}
\label{tab:circle-phase}
\begin{tabular}{lrr}
\toprule
Phase & Time (ms) & \% of Total \\
\midrule
Trace generation (GPU) & 0.3 & 1\% \\
NTT + trace commit (GPU) & 0.9 & 4\% \\
Fiat-Shamir & $<$0.1 & $<$1\% \\
Constraint eval + comp commit (GPU) & 2.0 & 10\% \\
FRI folding (CPU) & 12.0 & 57\% \\
Query phase & 6.0 & 28\% \\
\midrule
\textbf{Total} & \textbf{21} & 100\% \\
\bottomrule
\end{tabular}
\end{table}

FRI folding dominates at 57\% of total time, despite being executed on the
CPU.  The GPU phases (trace generation, NTT, constraint evaluation, Merkle
commitment) collectively account for only 15\% of proving time.  This
imbalance reflects two factors: (1)~the GPU phases benefit from
command buffer fusion (NTT + Merkle in one CB, constraint eval +
composition Merkle in one CB), reducing scheduling overhead to two command
buffer commits total; and (2)~CPU FRI folding involves $\log_2(Bn)$
sequential rounds of halving, each requiring a Merkle commitment and
transcript squeeze, which serializes naturally.

GPU-accelerating the FRI fold phase is a clear optimization target.  \sys
includes a \texttt{CircleFRIEngine} with GPU fold kernels (first fold via
$y$-twiddle, subsequent folds via $x$-twiddle, and a fused 2-round fold),
but the proving pipeline currently uses the CPU path for the full FRI
protocol to simplify Fiat-Shamir transcript management.  Integrating
GPU FRI folding with interleaved transcript squeezes would require
per-round GPU$\to$CPU synchronization (one command buffer commit per fold
round), which at ${\sim}$0.5\,ms per round would add ${\sim}$8\,ms of
dispatch overhead across 16~rounds---potentially negating the compute
savings at this trace size.

\paragraph{Proof sizes.}
Proof sizes range from 39\,KB ($2^8$) to 87\,KB ($2^{14}$), growing
sublinearly with trace size.  The dominant contribution is Merkle
authentication paths: each of the 16~queries opens $C = 2$ trace columns
and 1~composition polynomial, with path length $\log_2(Bn)$ hashes of
32~bytes each.  At $2^{14}$ with $4\times$ blowup, each path has 16~nodes
($2^{16}$ leaves), contributing $16 \times 3 \times 16 \times 32 =
24{,}576$~bytes per query set.  FRI round commitments and query
responses add approximately 20\,KB.

\subsubsection{Comparison with 256-bit Field Proof Systems}
\label{sec:circle-vs-plonk}

The Circle STARK's use of M31 arithmetic provides a per-operation
advantage over BN254-based systems, but STARKs inherently require more
operations than SNARKs: larger blowup factors (typically $4$--$16\times$
vs.\ none for Plonk/Groth16), more FRI rounds ($\log_2 n$ sequential
folds), and larger proofs (tens of KB vs.\ hundreds of bytes).
Table~\ref{tab:circle-vs-plonk} compares Circle STARK and Plonk at
roughly equivalent security and circuit complexity.

\begin{table}[t]
\centering
\small
\caption{Circle STARK vs.\ Plonk (BN254/KZG) on M3~Pro.}
\label{tab:circle-vs-plonk}
\begin{tabular}{lrr}
\toprule
 & Circle STARK & Plonk (KZG) \\
\midrule
Field & M31 (31-bit) & BN254 Fr (254-bit) \\
Prove (${\sim}$1K constraints) & 24\,ms & 49\,ms \\
Verify & 14\,ms & 2\,ms \\
Proof size & 53\,KB & $<$1\,KB \\
Field mul cost & 1 MAD & 64 MADs \\
FRI/PCS rounds & $\log_2 n$ & 1 (KZG) \\
Trusted setup & None & Required \\
\bottomrule
\end{tabular}
\end{table}

At roughly 1K constraints, Circle STARK proves $2\times$ faster than
Plonk---the M31 per-operation advantage ($64\times$) more than compensates
for the STARK's additional operations.  However, verification and proof
size strongly favor Plonk: KZG verification is a constant-time pairing
check (2\,ms), while STARK verification requires re-executing all query
checks and FRI verification (14\,ms).  Proof size differs by two orders of
magnitude (53\,KB vs.\ $<$1\,KB).

The crossover point depends on the application context.  For client-side
proving where proof generation latency matters most (mobile wallets,
browser provers), Circle STARK's faster proving and transparent setup are
advantages.  For on-chain verification where proof size and verification
cost dominate (ZK rollups), KZG-based systems remain preferable.

\subsubsection{Lessons for Metal GPU STARK Design}
\label{sec:circle-lessons}

The Circle STARK implementation yields several insights specific to the
Metal GPU architecture:

\paragraph{32-bit fields eliminate the ALU bottleneck.}
M31 field multiplication maps to a single 32-bit MAD instruction---Metal's
native throughput unit.  This transforms Metal's 32-bit ALU limitation
into a strength: the GPU's full datapath is utilized, unlike 256-bit
fields where $64\times$ more MAD instructions are needed per multiply.
The NTT + trace commit phase at $2^{14}$ completes in 0.9\,ms
(including two INTT, two NTT, and two Merkle trees), demonstrating that
M31 arithmetic is fast enough to make the NTT a negligible fraction of
total proving time.

\paragraph{Dispatch overhead, not compute, is the bottleneck.}
At all tested trace sizes ($2^8$ through $2^{14}$), GPU compute time for
the parallelizable phases (NTT, constraint evaluation, Merkle trees) is
under 4\,ms.  The proving time is instead dominated by sequential protocol
phases (FRI, query extraction) and Metal dispatch overhead.  This contrasts
sharply with BN254 workloads, where the 256-bit arithmetic itself dominates.
For M31 STARKs on Metal, future optimization effort should focus on
reducing the number of command buffer commits rather than accelerating
individual kernels.

\paragraph{Command buffer fusion is essential.}
\sys fuses logically distinct phases into shared command buffers: NTT +
Merkle commit in one CB, constraint evaluation + composition commit in a
second CB.  Without this fusion, the prover would require at least 6
command buffer commits (${\sim}$3\,ms of pure scheduling overhead),
increasing total time by ${\sim}$15\% at $2^{14}$.  At smaller trace
sizes, the impact is proportionally larger---explaining part of the
$2^8$ anomaly.

\paragraph{CPU FRI is a deliberate trade-off.}
Although \sys includes GPU Circle FRI fold kernels, the prover uses CPU FRI
to avoid per-round GPU$\to$CPU synchronization for Fiat-Shamir transcript
management.  At the current trace sizes ($\leq 2^{14}$), each FRI round
processes at most $2^{16}$ elements---small enough that CPU throughput
is adequate and the ${\sim}$0.5\,ms per-round dispatch overhead of a GPU
path would likely negate compute savings.  For larger traces ($2^{20}$+),
the balance shifts, and GPU FRI with batched fold rounds becomes
worthwhile.

\paragraph{The small-field thesis holds, with caveats.}
The migration from 256-bit fields to 31-bit fields delivers a clear
per-operation speedup on Metal: M31 multiply is $64\times$ cheaper than
BN254 Fr multiply, and the NTT benefits from $8\times$ better cache
utilization (4-byte vs.\ 32-byte elements).  However, the STARK protocol
structure (blowup, FRI rounds, larger proofs) partially offsets this
advantage.  The net effect at current trace sizes is a ${\sim}$2$\times$
proving speedup over Plonk at comparable constraint counts, with
significantly worse verification time and proof size.  The advantage
grows at larger trace sizes where GPU parallelism can amortize the
sequential FRI overhead.
